% !TEX root = ../AnonymousSubmission2027.tex
\section{Introduction}

% Contact-rich dexterous manipulation requires robots to continuously reason about local physical interactions, of which fingers are in contact, whether the object is slipping, and whether the multi-finger interaction remains stable. Currently, the widely used Vision-Language-Action models (VLAs) provide strong capabilities for instruction understanding and scene-level perception, but visual observations do not reliably reveal contact locations, force distributions, or grasp stability. This limitation becomes more severe during multi-finger manipulation, where the hand and object frequently occlude the contact regions. Tactile sensing directly measures local physical interactions and provides an essential feedback for fine-grained action
% adjustment and provides the feedback required for precise action adjustment.

Contact-rich dexterous manipulation requires robots to continuously reason about local physical interactions, such as which fingers are in contact, whether the object is slipping, and whether multi-finger contact remains stable. 
Vision-Language-Action models (VLAs) provide strong capabilities for instruction understanding and scene-level perception, but visual observations do not reliably reveal contact locations, force distributions, or grasp stability. 
This limitation is more pronounced in multi-finger manipulation, where the hand and object frequently occlude contact regions. 
Tactile sensing directly measures local physical interactions and provides essential feedback for fine-grained action adjustment.


% Recent work has integrated tactile or force sensing VLA policies in reactive and predictive manners. Reactive approaches~\cite{xue2025reactive, li2026favla, niu2026t} treat touch signals as additional observations alongside vision and robot states. They incorporate high-frequency tactile feedback into the control loop so the robot can quickly refine its actions when slip, contact misalignment, or force variations occur.
% Predictive approaches~\cite{zang2026tacforesight,niu2026learning,zhou2026touchworld}, in contrast, model future tactile signals, forces, or contact states and use the predicted interaction as a prior for subsequent action generation.
% Reactive methods enable rapid responses to observed contact changes, while predictive methods provide contact foresight about how the interaction may evolve.

Recent tactile-aware policies have progressively exploited touch beyond static observation. 
Direct-fusion methods incorporate current tactile signals as an additional input modality to improve contact awareness~\cite{huang2025tactile}. 
Reactive methods go further by using high-frequency tactile or force feedback to revise actions online in response to slip, contact misalignment, or force variations~\cite{xue2025reactive,li2026favla,niu2026t}. 
More recently, predictive methods anticipate future tactile signals or forces to guide action generation~\cite{zang2026tacforesight,niu2026learning}, while predictive--reactive approaches additionally combine such foresight with online action correction~\cite{zhou2026touchworld}. 
Together, these developments extend touch from an additional observation to reactive feedback and anticipatory control.

% Despite this progress, directly applying existing tactile-aware policies to dexterous manipulation remains challenging for two reasons. First, generic tactile tokenization is poorly suited to the dense and spatially structured tactile arrays of dexterous hands. As shown in Fig.~\ref{**}, each fingertip is equipped with a high-resolution tactile array and produces dense 3D force measurements. Encoding individual taxels as flattened tactile tokens~\cite{huang2025tactile,niu2026t} fails to preserve finger identity and local contact regions of the fingertip surface. 
% As a result, the model may detect a change in the overall tactile state but struggle to determine which finger requires adjustment, limiting region-level tactile perception and precise finger-wise control. 
% Second, predicting future contact is challenging for dexterous hands, whose multi-finger motions have higher degrees of freedom than two-finger grippers. Existing methods typically generate a single future tactile prediction and then use it to produce an open-loop action chunk. As a result, the action is driven by a potentially inaccurate contact forecast, without using high-frequency tactile feedback to timely refine the action and forecast jointly.

% 尽管取得了这些进展，现有触觉感知策略仍不足以应对接触丰富的灵巧操作，主要存在两个原因。首先，灵巧手触觉同时具有密集采样、局部稀疏响应和明确空间拓扑。将单个触觉单元展平为通用触觉 token~\cite{huang2025tactile}，无法显式保留手指身份以及各指尖的功能性接触拓扑，使得策略难以将触觉变化准确关联到需要调整的具体手指和接触区域。因此，灵巧操作需要一种能够按照手指身份和局部接触区域组织原始触觉信号的结构化表示。
% 其次，多指运动会产生复杂且快速变化的接触动力学，使灵巧手的未来接触比二指抓夹更难被准确预测。现有预测式方法通常只在动作片段执行前生成触觉预见，并在执行周期中保持固定。然而，滑移、接触错位和执行误差都可能迅速导致该触觉预见与当前交互状态不一致。尽管新观测到的触觉反馈可以用于修正动作，但指导该动作的触觉预见并未得到同步更新，由此产生预测与控制之间的不一致。
Despite this progress, existing tactile-aware policies remain insufficient for contact-rich dexterous manipulation for two reasons. 
First, dexterous-hand touch combines dense sampling, locally sparse responses, and explicit spatial topology. 
Flattening individual taxels into generic tactile tokens~\cite{huang2025tactile} does not explicitly preserve finger identity or the functional topology of fingertip contact regions, making it difficult to associate tactile changes with the specific finger and region that require adjustment. 
Dexterous manipulation therefore calls for a structured representation that organizes raw tactile signals by finger identity and local contact region.
Second, tightly coupled multi-finger motions induce complex and rapidly changing contact dynamics, making future contact harder to predict accurately than with two-finger grippers. 
Existing predictive methods typically generate tactile foresight before executing an action chunk and keep it fixed throughout the execution cycle. 
Slip, contact displacement, or execution errors can quickly make this foresight inconsistent with the current interaction state. 
Although newly observed tactile feedback may revise the action, the foresight guiding that action is not updated accordingly, resulting in a prediction--control mismatch.



% 为解决上述问题，我们提出 \systemname，一种面向接触丰富灵巧操作、具有递归触觉预见能力的触觉引导 VLA 框架。其 Tactile Patch Encoder 按照手指身份和仿人指尖拓扑组织原始触觉信号，将每根手指的 taxels 聚合为结构化 patch features，再通过接触感知的注意力池化形成每根手指一个 tactile token，从而以固定的 token 预算保留局部接触结构。基于这些结构化触觉表示，我们进一步学习 action-relevant future tactile latents。两个动作专家共享动作生成监督：Hindsight Action Expert 利用真实未来触觉提取 latent targets，Foresight Action Expert 则根据当前可用观测预测相应的 latents，并通过 latent alignment 与前者对齐。
% 我们进一步提出异步闭环执行机制，递归修正触觉预见与动作。在每个 action chunk 中，VLM 仅编码一次视觉语言上下文并缓存结果，Foresight Action Expert 则被高频调用。每次调用结合最新触觉历史，用上一轮预测的 tactile latent prefix 初始化当前 foresight queries 的对应前缀，随后重新预测 future tactile latents 并更新动作。非对称因果注意力使 foresight queries 无法访问 action tokens，同时允许动作读取更新后的 tactile latents，使触觉预见单向引导动作生成。由此，触觉预见从一次性预测转变为随物理交互递归维护的控制表示。
% To address these limitations, we introduce \systemname, a closed-loop VLA model with recursive tactile foresight for contact-rich dexterous manipulation.
% We employ a decoupled VLA architecture with the low-frequency VLM reasoning and high-frequency tactile closed-loop action expert. 
% First, we introduce a \Encoder  to\textcolor{blue}{need to fix} extract spatially structured representations from dense multi-finger tactile arrays. This encoder organizes raw taxels into structured patch features, encoded by local contact regions, fingertip positions, and finger identities. This representation provides region-level tactile perception and spatially grounds action adjustments to the corresponding fingers and contact areas.
% Second, with the structured tactile representations, we present an action-relevant tactile foresight learning framework during training with dual action experts. Specifically, A hindsight expert leverages privileged future tactile inputs to produce future latents, while a causal foresight expert predicts these latents from current observations alone.  Instead of reconstructing future tactile maps, we align predicted and future latents with a latent alignment loss. A shared action flow matching loss across both experts extracts action-relevant latents and preserves information about future multi-finger interactions.
% Finally, we propose an asynchronous closed-Loop tactile foresight inference strategy to recursively refine the action and future latent jointly. \systemname\ caches the low-frequency semantic context from the VLM backbone and runs a high-frequency foresight action expert multiple times per action chunk. \textcolor{blue}{check:}Each iteration conditions on the latest tactile context and previously predicted future latents to refine the unobserved future latents and remaining actions. This keeps tactile foresight and fine-grained control synchronized as new contact observation arrives.

To address these limitations, we introduce \systemname, a contact-guided VLA framework with recursive tactile foresight for contact-rich dexterous manipulation. We develop a Tactile Patch Encoder that organizes taxels by finger identity and human-inspired fingertip topology. It aggregates each finger's taxels into structured patch features and uses contact-aware attention pooling to produce one token per finger, preserving local contact structure within a fixed token budget. Based on these representations, dual action experts learn action-relevant future tactile latents under shared action-generation supervision. The Hindsight Action Expert extracts latent targets from ground-truth future touch, while the Foresight Action Expert predicts and aligns the corresponding latents from currently available observations. We further introduce asynchronous closed-loop execution to recursively revise tactile foresight and actions. Within each action chunk, the VLM encodes and caches the visual-language context once, while the Foresight Action Expert runs at a higher frequency. Each invocation uses the latest tactile history and the previous tactile latent prefix to initialize the current foresight-query prefix, then updates the future tactile latents and remaining actions. An asymmetric causal mask prevents foresight queries from attending to action tokens while allowing actions to attend to updated tactile latents, so tactile foresight directly guides action generation. Tactile foresight thus becomes a control representation recursively maintained throughout interaction.


% 为评估 \systemname，我们搭建了 XHand--UR7e 真实触觉操作平台，并构建了一个包含 800 条真机示范和七个接触丰富任务的灵巧触觉操作基准。在该基准上，\systemname 的平均成功率相比最强基线提升了 [X] 个百分点。在涵盖光照、背景、物体高度和物体颜色变化的分布外测试中，\systemname 同样表现出更强的泛化能力，平均成功率相比最强基线提升了 [Y] 个百分点。消融实验进一步验证了 \Encoder、双动作专家设计以及触觉预见与动作异步联合修正机制的有效性。
To evaluate \systemname, we build a real-world XHand--UR7e tactile manipulation platform and construct the XHand Tactile Benchmark (XHT-Bench), comprising 800 real-world demonstrations across seven contact-rich tasks. On this benchmark, \systemname improves the average success rate over the strongest baseline by [X] percentage points. In out-of-distribution tests covering variations in lighting, background, object height, and object color, \systemname also exhibits stronger generalization, outperforming the strongest baseline by [Y] percentage points in average success rate. Ablation studies further validate the effectiveness of the \Encoder, the dual-action-expert design, and the asynchronous co-revision of tactile foresight and actions.

Our contributions are summarized as follows:

% 我们提出 \systemname，一种具有递归触觉预见能力的接触引导 VLA 框架。它利用最新触觉反馈异步修正 future tactile latents 与剩余动作，闭合触觉预见与动作控制之间的反馈循环。

% 我们提出面向灵巧手的结构化触觉表示学习方法。Tactile Patch Encoder 按照手指身份和指尖接触拓扑组织密集触觉信号，并通过双动作专家在动作生成监督下学习 action-relevant future tactile latents。

% 我们搭建了 XHand--UR7e 触觉操作平台并构建 XHT-Bench，涵盖 800 条真机示范和七个接触丰富的任务。\systemname 在标准和分布外测试中的平均成功率分别领先最强基线 [X] 和 [Y] 个百分点，消融实验进一步验证了关键设计。
\begin{itemize}
    \item We propose \systemname, a contact-guided VLA framework with recursive tactile foresight. Using the latest tactile feedback, \systemname\ asynchronously revises future tactile latents and remaining actions, closing the feedback loop between tactile foresight and action control.

    \item We introduce a structured tactile representation learning method for dexterous hands. The \Encoder organizes dense tactile signals by finger identity and fingertip contact topology, while dual action experts learn action-relevant future tactile latents under action-generation supervision.

    \item We build an XHand--UR7e tactile manipulation platform and construct XHT-Bench, comprising 800 real-world demonstrations across seven contact-rich tasks. \systemname\ outperforms the strongest baseline by [X] and [Y] percentage points in average success rate under standard and out-of-distribution settings, respectively, while ablations validate its key components.
\end{itemize}




% 完成接触密集型灵巧操作，需要机器人在执行过程中持续判断接触是否建立、物体是否滑移以及交互是否稳定。视觉-语言-动作模型（Vision-Language-Action models, VLAs）擅长指令理解和场景感知，但视觉难以可靠揭示接触位置、受力分布和抓持稳定性等局部物理状态，多指操作中的严重遮挡进一步加剧了这一局限。触觉能够直接感知局部接触，为精细动作调节提供关键反馈，因此是实现稳健灵巧操作的重要感知模态。
% Contact-rich dexterous manipulation requires robots to continually assess whether contact has been established, whether an object is slipping, and whether the interaction remains stable. Vision-language-action models (VLAs) can interpret instructions and perceive scene structure, but vision alone does not reliably reveal local physical states such as contact location, force distribution, and grasp stability. Severe occlusion during multi-finger manipulation further compounds this limitation. Tactile sensing directly captures local contact and provides the feedback needed for fine-grained action adjustment and robust dexterous manipulation.




% 现有方法对触觉的利用大致可分为反应式与预测式两类。反应式方法或者将触觉作为额外观测，与视觉和state共同用于策略学习；或者将高频触觉反馈引入执行环路，在线修正动作以应对滑移、接触错位和受力变化。这些方法增强了策略对接触变化的适应能力，但其决策主要建立在已经观测到的触觉结果之上，通常只能在接触偏差出现后进行调整。为进一步获得对接触演化的前瞻性，近期的预测式方法开始建模未来触觉、力或接触状态，并将其作为动作生成的预测条件或执行参考。
% Existing work broadly falls into two categories based on how it uses tactile information: reactive and predictive. Reactive approaches either treat touch as an additional observation alongside vision and robot state, or incorporate high-frequency tactile feedback into the execution loop to revise actions online in response to slip, contact misalignment, and force variation. These approaches improve adaptation to changing contact conditions. However, their decisions depend primarily on tactile outcomes that have already been observed, so correction typically occurs only after contact deviates from the intended interaction. To anticipate how contact will evolve, recent predictive approaches model future tactile signals, forces, or contact states and use them as conditions for action generation or references during execution.

% 然而在现有预测式或预测－反应式方法中，未来触觉通常在规划周期开始时就生成了；在随后的执行过程中，新到来的触觉反馈主要用于修正动作，而未来触觉表示本身通常不会随之递归更新。这忽略了未来触觉由当前接触状态与后续动作共同决定：一旦滑移、物体形变、接触偏移或执行误差改变了交互状态，修正后的动作也将进一步改变后续接触。此时，原有预测便不再对应当前状态和新的动作计划。若策略只更新动作，却继续依赖对应于原始计划的未来触觉表示，已观测触觉、未来接触预期与后续动作便会逐渐失去一致性。因此，真正缺失的不只是更准确的未来预测或更快速的动作修正，而是将未来触觉预见本身纳入反馈回路，根据最新触觉证据联合修正未来触觉表示和尚未执行的动作。
% However, existing predictive and predictive-reactive approaches typically generate a future tactile representation once at the beginning of a planning cycle. During execution, incoming tactile feedback is used primarily to revise actions, while the forecast itself is left fixed rather than updated recursively. This design overlooks that future touch is jointly determined by the current contact state and the actions yet to be executed. Slip, object deformation, contact shifts, or execution errors can change the interaction state. The resulting action revisions further alter how contact will evolve, leaving the original forecast mismatched with both the current state and the revised action plan. If the policy updates only its actions while retaining a tactile forecast tied to the original plan, the observed tactile history, expected future contact, and remaining actions become increasingly inconsistent. The key missing component is a feedback loop that uses the latest tactile evidence to jointly revise the future tactile representation and the unexecuted actions. Improving prediction accuracy or action correction speed alone does not provide this coupling.


% % 为解决上述问题，我们提出 \systemname，一种面向灵巧操作、具有闭环触觉预见能力的 VLA。\systemname\ 会高频联合修正 future tactile latents 与未执行动作，从而将低频视觉语言推理与高频触觉闭环控制解耦。具体来讲，在执行过程中，模型缓存视觉语言模块提取的语义上下文，并在相邻两次语义更新之间多次调用高频触觉动作分支。每次调用均以已观测触觉上下文和上一轮预测的未来触觉状态为条件，递归更新 future tactile latents，并更新动作块。因此，每轮动作修正都以当前轮更新的触觉预见为条件：该预见既继承上一轮预测中的未来信息，又融合执行过程中持续获得的最新触觉证据，从而使触觉预见能够随当前交互状态和动作计划共同调整，减少持续执行中的失配，避免后续控制持续依赖过时的一次性预测。为有效表征灵巧手上高维、分散且结构复杂的触觉信号，我们设计了 Tactile-Patch Encoder，按照接触拓扑将原始触觉组织为结构化的 tactile patch features。在此基础上，我们利用动作生成监督进一步提炼真实未来触觉中的控制相关信息，形成紧凑的 action-relevant tactile representations。【策略预测的 future tactile latents 与这些表征进行对齐，从而学习对后续动作生成有用的触觉预见。】
% To address this limitation, we introduce \systemname, a VLA with closed-loop tactile foresight for dexterous manipulation. \systemname\ jointly revises future tactile latents and the remaining unexecuted actions at high frequency. To support this process, it decouples low-frequency vision-language reasoning from high-frequency tactile closed-loop control. During execution, the model caches the semantic context extracted by the vision-language module and invokes a high-frequency tactile-action branch multiple times between consecutive semantic updates. Each invocation conditions on the observed tactile context and the future tactile state predicted in the previous iteration. It then recursively updates the future tactile latents and revises the remaining actions in the current action chunk. Each action revision is therefore conditioned on the tactile foresight updated in the same iteration. This foresight retains future information from the previous prediction while incorporating the latest tactile evidence gathered during execution. It can thus evolve with the current interaction state and the revised action plan, reducing prediction-control mismatch and preventing the policy from relying on a stale one-shot forecast. To represent the high-dimensional, spatially distributed, and structurally complex tactile signals of a dexterous hand, we design a Tactile-Patch Encoder that organizes raw tactile observations into structured tactile patch features according to the contact topology. During training, action-generation supervision further extracts control-relevant information from ground-truth future tactile observations, producing compact action-relevant tactile representations.









% 我们的贡献概括如下：


